% Generated draft: requires mathematical and source review.
\documentclass[11pt]{article}
\usepackage[margin=25mm]{geometry}
\usepackage{fontspec}
\setmainfont{Latin Modern Roman}
\usepackage{amsmath,amssymb,mathtools,unicode-math}
\setmathfont{Latin Modern Math}
\usepackage{enumitem,needspace,xurl,hyperref}
\setlength{\parindent}{0pt}
\setlength{\parskip}{5pt}
\setlength{\emergencystretch}{3em}
\newcommand{\N}{\mathbb{N}}
\newcommand{\Z}{\mathbb{Z}}
\newcommand{\Q}{\mathbb{Q}}
\newcommand{\R}{\mathbb{R}}
\newcommand{\C}{\mathbb{C}}
\newcommand{\F}{\mathbb{F}}
\newcommand{\A}{\mathbb{A}}
\newcommand{\PP}{\mathbb{P}}
\newcommand{\E}{\mathbb{E}}
\newcommand{\eps}{\varepsilon}
\newcommand{\dd}{\,\mathrm{d}}
\begin{document}
\section*{UnsolvedMath ID: 30001913}
\textbf{Source problem code:} catalogue entry
\subsection{Definitions and mathematical statement}
A Laurent polynomial in finitely many variables has the form
\[
f(x_1,\ldots,x_n)=\sum_{m\in S}c_m x_1^{m_1}\cdots x_n^{m_n},\qquad S\subseteq\Z^n\text{ finite}.
\]
The supplied record does not specify the coefficient domain, the allowed values of \(n\), or additional admissibility hypotheses.

\textbf{Conjecture, with unresolved source-dependent terminology.} Every Laurent polynomial that is Hodge--Tate at infinity is extremal. Equivalently, writing \(\mathcal L\) for the class of Laurent polynomials intended in the cited report, the asserted implication is
\[
\forall f\in\mathcal L,\qquad
\bigl(f\text{ is Hodge--Tate at infinity}\bigr)
\Longrightarrow
\bigl(f\text{ is extremal}\bigr).
\]
This is a transcription of the supplied conjecture, not a complete intrinsic formulation: the record does not define \(\mathcal L\), identify the geometric or cohomological construction used to define ``Hodge--Tate at infinity,'' or define ``extremal.'' These predicates must retain their intended source-specific meanings; no replacement definitions or additional hypotheses are asserted here.
\subsection{Short English statement}
The conjecture asks whether the condition called Hodge--Tate at infinity forces a Laurent polynomial to be extremal. The supplied record omits the definitions and admissibility hypotheses needed to make this implication fully precise. Its open status is supplied dated curation, not an independently verified conclusion.
\subsection{Sources}
\begin{itemize}
\item Local registry: supplied UnsolvedMath record, permanent ID \(30001913\), title ``Extremality of Hodge--Tate Laurent Polynomials,'' exact source problem code \texttt{catalogue entry}, in the Oberwolfach Reports Open Problems collection. The record assigns proposed year \(2011\) and identifies no individual proposer.
\item Supplied primary citation: \emph{Complex Algebraic Geometry (2012)}, Oberwolfach Reports, pp.~2669--2733, DOI \texttt{10.4171/owr/2011/47}. \url{https://doi.org/10.4171/owr/2011/47}. The supplied background describes these reports as documenting mathematical workshops and their open-problem discussions. No pinpoint location for this conjecture is supplied.
\item Supplied literature reference: Coates--Kasprzyk--Pitton--Tveiten, ``maximally mutable Laurent polynomials.'' \url{https://arxiv.org/abs/2107.14253}.
\item The supplied literature-source list also identifies the original OWR report: \url{https://doi.org/10.4171/owr/2011/47}.
\item Literature assessment dated exactly \texttt{2026-08-21}: the supplied curation labels the problem open, says that the conjecture is explicitly stated, and reports that the located later work does not prove it in full generality. This is supplied dated literature triage, not an independently verified assessment; the record requests independent verification of time-sensitive status and cited claims.
\item The registry labels the statement \texttt{corrected\_verified} and reports that its mathematical content was checked against the official report, with notation, context, and question form normalized without changing the source-level claim. This verification claim is supplied provenance and has not been independently checked for this draft.
\end{itemize}
\end{document}
